(i) In $\Cl(V) \otimes \C$, for an orthonormal basis put $\epsilon_i = e_i$ if $Q(e_i) = 1$ and $\epsilon_i = e_i \otimes I_4$ if $Q(e_i) = -1$. The $\epsilon_i$ anticommute and square to $1$, so the complexified universal property gives a homomorphism $\Cl_n(\C) \to \Cl(V) \otimes \C$, onto because the $\epsilon_i$ and the scalar $1 \otimes I_4$ generate, and an isomorphism since both have complex dimension $2^{n}$. So $\Cl_n(\C) \cong \Cl(n, 0) \otimes \C$; the complexified \claimref{std-algebra.tensor-rules} with $W = \R^{1,1}$ give $\Cl_{n+2}(\C) \cong \Mat_2(\Cl_n(\C))$, and $\Cl_0(\C) = \C$, $\Cl_1(\C) = (\R \oplus \R) \otimes \C = \C \oplus \C$ (\claimref{std-algebra.periodicity}) start the induction.

(ii) For odd $n$, $\omega$ is central (\claimref{primer.centre}) and $\omega^{2} = \pm 1$ (\claimref{primer.sign}), so $\omega$ is invertible. It is odd, so multiplication by $\omega$ maps $\Cl^{+}$ injectively into $\Cl^{-}$, which has the same dimension $2^{n-1}$ (\claimref{std-algebra.z2-grading}); hence $\Cl^{-} = \Cl^{+}\omega$ and $\Cl = \Cl^{+} \oplus \Cl^{+}\omega$. If $\omega^{2} = 1$, $P_{\pm}$ are central with $P_\pm^{2} = P_\pm$, $P_+ P_- = 0$, $P_+ + P_- = 1$, so $\Cl = \Cl P_+ \oplus \Cl P_-$ as algebras. The map $a \mapsto aP_+$ is a homomorphism on $\Cl^{+}$ because $P_+$ is a central idempotent; it is injective because $aP_+ = \tfrac12(a + a\omega) = 0$ forces the even part $a$ to vanish; $\Cl P_+ = \Cl^{+} P_+$ since $\omega P_+ = P_+$; likewise for $P_-$. If $\omega^{2} = -1$, the map is a homomorphism because $\omega$ is central and squares to $-1$, and a bijection by $\Cl = \Cl^{+} \oplus \Cl^{+}\omega$.
